% P5 method subsection (draft, 20 Oct 2026). Budget: about 0.4 page of the 6-page AAAI report.
% Cite key park2025decoding = Park, Kim & Ko, "Decoding Dense Embeddings: Sparse Autoencoders for
% Interpreting and Discretizing Dense Retrieval", EMNLP 2025, pp. 26468--26485.
% Numbers marked \todo{} wait for the 3 Nov freeze.

\subsection{Query Concept Expansion (P5)}
\label{sec:p5}

CL-SR scores a passage only through concepts it shares with the query, and a short query activates few of
them. A relevant passage written with related concepts can therefore score zero, which is the
vocabulary-mismatch failure that \citet{park2025decoding} measure on the 988 MS MARCO dev queries where BM25
misses every relevant passage in its top 1000; there CL-SR Max trails SimLM by 23\% MRR@10 (0.143 against
0.185). We change only the query side: passages, the index, the SAE and the BM25 constants stay as released,
and the original query concepts keep their weights.

\paragraph{Candidate pool.}
For each query we collect two kinds of candidate concepts. \emph{Near-miss} concepts are the 32 with the
largest positive SAE pre-activation that still fall at or below the inference threshold; the encoder already
half-detects them. \emph{Decoder neighbours} are, for each active concept, the 8 concepts whose decoder
directions have the highest cosine to it. We drop concepts already in the query, concepts absent from the
index, and concepts with $\mathrm{idf}\le 0$, since the released BM25 IDF is negative for concepts in more
than half of the passages and adding them can only lower scores.

\paragraph{Expander.}
A two-layer MLP with 1.3k parameters maps five features of each pool concept $c$ (its pre-activation, the
cosine between its decoder direction and the SimLM query vector, its largest weighted cosine to an active
concept, its IDF, and the log number of active concepts) to a weight $w_c=\mathrm{softplus}(\cdot)\ge 0$.
Scoring pool candidates one at a time keeps the model small and means it cannot add a concept outside the
filtered pool.

\paragraph{Training through the released scorer.}
The student score of candidate passage $d$ is the CL-SR formula itself,
\begin{equation}
\begin{split}
  s(q',d)=\sum_{c\in q'\cap d}\mathrm{idf}(c)\,
  &\frac{(k_2+1)\,w_c}{k_2+w_c}\\
  \cdot\,&\frac{(k_1+1)\,f_{d,c}}{f_{d,c}+k_1\bigl(1-b+b\,|d|/\overline{|d|}\bigr)},
\end{split}
\end{equation}
with the constants of the released code ($k_1=0.6$, $k_2=1.5$, $b=2.0$), BM25 IDF and $|d|$ the number of
indexed concepts; the code and the paper's Eq.~5 differ here, and we follow the code. Only the query
saturation term depends on the expander, so the passage factors are precomputed. We distil SimLM over each
training query's dense top-100 candidates,
\begin{equation}
\begin{split}
  \mathcal{L}=T^2\,\mathrm{KL}\bigl(&\mathrm{softmax}(s_{\text{SimLM}}/T)\,\big\|\\
  &\mathrm{softmax}(\alpha\, s(q',\cdot)/T)\bigr)+\lambda\sum_{c} w_c ,
\end{split}
\end{equation}
where $\alpha$ is a learned scale and the L1 term pushes most pool weights towards zero. Training the
expander on the scorer it will be evaluated with removes any proxy gap; a concept-space dot product would
ignore the saturation, IDF and length terms that decide CL-SR rankings.

\paragraph{Budget.}
At inference the expander keeps its top-$B$ concepts, which bounds the extra postings each query reads. We
choose $B\in\{4,8,16\}$ and the checkpoint jointly by MRR@10 when re-ranking the validation top-100, and
report final numbers from full-index retrieval.

\paragraph{Controls.}
Because near-miss concepts make part of the expander a per-query threshold, we compare against a
\emph{query-threshold} row (T) that re-encodes queries with the threshold scaled by $s\in\{0.9,\dots,0.5\}$,
tuned on validation, and against a rule-based \emph{decoder-neighbour} row (N) that adds the top-$B$
neighbours with weight equal to the source concept's weight times their cosine. Significance comes from a
paired randomisation test (10{,}000 sign flips) on the full query set and on the BM25-mismatch subsets, which
we rebuild with Pyserini BM25 so that validation queries have them too. Our claims rest on six tests, L
against B0, T and N on Recall@1000 and on mismatch MRR@10, which we correct together with Holm's method at
$\alpha=0.05$; other p-values in the table are uncorrected. Every query counts in the denominator, including those for which CL-SR returns nothing; the
released evaluation script skips them, which would favour expansion.
